%!TEX root = main.tex
\paragraph{Topology versus price shocks: a quantified gap.}
The scenario library spans a wide range of disruption types, but
the results reveal a clean hierarchy. Topological shocks,
those that physically disconnect suppliers from the network,
are categorically more destructive than price shocks of any
severity. The worst purely price-based scenario, T3 (hub
surcharge $+30\%$), costs $\Delta Z_F = +\$852{,}284$
($+19.1\%$). The worst topological scenario, S3 (H3 closure
combined with energy shock), costs
$\Delta Z_F = +\$3{,}989{,}134$ ($+89.6\%$). The ratio is:
\[
  \frac{\Delta Z_F(\text{S3})}{\Delta Z_F(\text{T3})}
  = \frac{89.6\%}{19.1\%} \approx 4.7\times
\]
Topological shocks are roughly five times more
destructive than the most severe price shock tested, even
before accounting for the unmet-demand dimension of S1 and S3
under strategies R and A. The intuition is structural: a price
shock raises variable costs and the solver can partially reroute
around the most expensive lanes; a topological shock severs
connectivity entirely, leaving no routing alternative and forcing
spot-market procurement at a $2.5\times$ premium. Scenario~S4
(Suez corridor, $+7.7\%$) confirms this: an $+80\%$ cost
increase on 39 sea arcs, severe by any measure, propagates
as only a $+7.7\%$ system-wide cost increase because alternative
routes exist. The H3 closure propagates as $+81.6\%$ because
they do not.
 
\paragraph{Modal exposure and product vulnerability.}
The GlobalFlow instance uses exclusively two transport modes:
sea (240 arcs, long-haul inter-node legs) and
road (77 arcs, last-mile warehouse-to-customer
delivery). No air freight arcs are present in the network.
Table~\ref{tab:modal_split} reports the resulting modal split
by product in the baseline optimal solution.
 
\begin{table}[H]
\centering
\caption{Baseline modal split by product — flow volume and
         transport cost share.}
\label{tab:modal_split}
\small
\begin{tabular}{@{}lcccc@{}}
\toprule
& \multicolumn{2}{c}{\textbf{Sea}} &
  \multicolumn{2}{c}{\textbf{Road}} \\
\cmidrule(lr){2-3}\cmidrule(lr){4-5}
\textbf{Product} & Volume & Cost & Volume & Cost \\
\midrule
A — Fertilizers        & 75\% & 80\% & 25\% & 20\% \\
B — Semiconductors     & 82\% & 96\% & 18\% &  4\% \\
C — Battery Components & 74\% & 91\% & 26\% &  9\% \\
\bottomrule
\end{tabular}
\end{table}
 
Three findings stand out. First, the absence of air freight has
a direct implication for scenario~T2 (energy crisis: sea~$+10\%$,
road~$+20\%$, air~$+60\%$): the $+60\%$ air surcharge is
structurally irrelevant for GlobalFlow, the network
carries no air flows to penalise. The effective shock is limited
to sea and road, which is why all three strategies tie exactly
under T2 ($Z_R = Z_A = Z_F$): the baseline modal mix is already
optimal under a uniform per-mode multiplier, and no rerouting or
reconfiguration can improve on it.
 
Second, product~B (Semiconductors) is the most
sea-intensive by cost (96\% of its variable transport cost),
making it the largest absolute contributor to the T2 disruption
cost ($+\$293{,}217$, or $+10.4\%$ on its variable cost base).
Product~A (Fertilizers), despite lower sea intensity by cost
(80\%), is proportionally \textit{more} affected ($+12.0\%$)
because its road component carries a higher cost weight relative
to product~B, and road is penalised at $+20\%$ versus sea's
$+10\%$.
 
Third, canal disruptions (S4) are structurally
contained by the road share. The 18--26\% road volumes of all
three products are entirely unaffected by a Suez corridor shock,
providing a natural partial buffer. This explains why even the
most severe canal scenario ($+7.7\%$) remains far below the
modal-level energy shock ($+10.6\%$): the road component acts
as a cost-insensitive cushion that limits the propagation of any
sea-corridor disruption.